% Page 089

\seite{89}

Wir gehen wieder am Beginn eines neuen\ob
Jahres. Unsere Lage hat sich im vergangenen Raum\ob
gebessert. Immer noch stehen fremde Truppen\ob
in unserer eigenen Heimat. Wie lange noch?\ob
Möchte es mit die so sehr und groß\unsicher Zahl der\ob
Arbeitslosen fort von dem Wohlstand, den uns\ob
unser Herrgott so gerne zuteilen! Wenn\ob
von den Feinden, so lebt man sich\unsicher nach den\ob
ersten Jahr Besserung die versprochenen Ver=\ob
billigung. Gebe es Gott, daß es einst wieder die\ob
frühere Vaterland noch wie den Einzelnen\ob
wieder ein Jahr die Enttäuschungen werde!

\pend
\abschnitt{In Gottes Namen fangen wir an!}
\pstart

\margmark{3. Januar    B}%
ei der gestrigen Eröffnung der Weihnachtsferien\ob
zeigte es sich, daß von den 44 Kindern noch 28 krank\ob
sind. Die Folge ist die erneute Schulschließung bis\ob
einschließlich 15.\ Januar.

Unsere Einquartiert am Jahresbeginn wurde\ob
für die Familie Großvater Schmitz Trauer. Eine 24 jäh=\ob
rige Tochter bei auf ihren Krankheit im Hotel\ob
Müller in Kyllburg in der Rückenmarkserkrankung am\ob
Rücken erkrankt. Noch die Familie wurde sich allge=\ob
mein Teilnahme zu.

\margmark{1. Februar   I}%
mmer noch sind mehrere Kinder an den Fol=\ob
gen der Masern so krank, Keuchhustenerscheinung.\ob
Mit Schulanfang von den 41 Schulkindern\ob
waren 37 an den Masern so krank. Besprochen\ob
wissen fast sämtliche nicht schulpflichtige Kinder.\ob
Das Dorf sieht dieses Krankheit behaftet. Im Alt=\ob
ter von 3--4 Jahren Schmitz\unsicher Jos. Joh. Nik.\unsicher sind daran ge=\ob
storben.
\pend
